% Part: incompleteness
% Chapter: introduction
% Section: overview

\documentclass[../../../include/open-logic-section]{subfiles}

\begin{document}

% SEGMENT OLP-0278-S01

\olfileid{inc}{int}{ovr}

\olsection{முழுமையின்மை முடிவுகளின் மேலோட்டம்}

கணிதத்தை !!{axiomatizable} கோட்பாடாக முறைப்படுத்த முடியும் என்றும், அந்தக்
கோட்பாடு !!{complete} ஆகவும் !!{decidable} ஆகவும் இருப்பதை நிறுவ முடியும்
என்றும் ஹில்பர்ட் எதிர்பார்த்தார். மேலும், உள்ளுணர்வியலின் அறைகூவல்களுக்கு
எதிராகச் செவ்வியல் கணிதத்தைக் காக்கும் வகையில், மிக வலுவற்ற,
``முடிவுறுமுறைச்'' சாதனங்களால் இக்கோட்பாட்டின் முரணின்மையை நிறுவுவதை அவர்
நோக்கமாகக் கொண்டார். இந்த நோக்கங்களை அடைய முடியாது என்பதை கோடலின்
முழுமையின்மைத் தேற்றங்கள் காட்டின.

ரசல் மற்றும் ஒயிட்ஹெட்டின் \emph{பிரின்சிபியா மேத்தமாட்டிக்கா}வின் ஒரு
வடிவம் !!{complete} அல்ல என்பதை கோடலின் முதல் முழுமையின்மைத் தேற்றம்
காட்டியது. ஆனால் அதன் நிறுவல் உண்மையில் மிகவும் பொதுவானது; பலவகையான
கோட்பாடுகளுக்கும் அது பொருந்துகிறது. இதன் பொருள்,
\emph{பிரின்சிபியா மேத்தமாட்டிக்கா} கணிதத்தை முழுமையாகப் பற்றிக்கொள்ளத்
தவறியது என்பது மட்டும் அல்ல; ஏற்றுக்கொள்ளத்தக்க \emph{எந்தக்} கோட்பாடும்
அவ்வாறு பற்றிக்கொள்ள முடியாது என்பதே. முழுமையின்மைத் தேற்றங்கள் பொருந்தப்
போதுமான கோட்பாட்டுப் பண்புகளைத் தனிமைப்படுத்தவும், கோடலின் நிறுவலை அந்தப்
பண்புகளை மட்டுமே சாருமாறு பொதுமைப்படுத்தவும் சிறிது காலம் எடுத்தது. ஆனால்
எண்கணிதத்தின் மொழியான~$\Lang L_A$ இலுள்ள கோட்பாடுகளுக்கான முதல்
முழுமையின்மைத் தேற்றத்தின் மிகவும் பொதுவான வடிவத்தை இப்போது கூறலாம்.

\begin{thm}
$\Gamma$ ஆனது~$\Lang L_A$ இலுள்ள முரணற்ற, !!{axiomatizable} கோட்பாடாகவும்,
எல்லாக் கணிக்கத்தக்க சார்புகளையும் !!{decidable} தொடர்புகளையும்
!!{represents} செய்வதாகவும் இருந்தால், $\Gamma$ !!{complete} அல்ல.
\end{thm}

$\Gamma$ !!{complete} அல்ல என்று சொல்வது, குறைந்தது ஒரு
!!{sentence}~$!A$ க்கு $\Gamma \Proves/ !A$ மற்றும் $\Gamma \Proves/
\lnot !A$ ஆகிய இரண்டும் நிலவுகின்றன என்று சொல்வதாகும். அத்தகைய
% SOURCE CORRECTION TA-INC-001: keep the parenthesis outside the inline formula.
!!a{sentence} ஆனது ($\Gamma$ வைச் சாராத) \emph{சார்பற்ற வாக்கியம்}
எனப்படுகிறது. சார்பற்ற வாக்கியங்கள் இருக்கவேண்டும் என்பதை உண்மையில் மிக
விரைவாக நிறுவ முடியும். ஆனால் அத்தகைய சார்பற்ற !!{sentence} க்கு ஒரு
\emph{குறிப்பிட்ட எடுத்துக்காட்டை} கோடலின் நிறுவல் தருவதில்தான் அதன் வலிமை
அடங்கியுள்ளது. அந்த ஆர்வமூட்டும் கட்டமைப்பு, $\Gamma$ வுக்கான
\emph{கோடல் வாக்கியம்} எனப்படும்~$!G_\Gamma$ என்ற !!a{sentence} ஐ
உருவாக்குகிறது. $\Gamma$ வில், $!G_\Gamma$ ஆனது $\Gamma$ விலிருந்து
$!G_\Gamma$ வருவிக்கத்தக்கதல்ல என்ற கூற்றுக்குச் சமானமாக இருப்பதால் அது
வருவிக்கத்தக்கதல்ல. இதை அது \emph{கட்டமைப்புமுறையில்} செய்கிறது: அதாவது,
$\Gamma$ வின் அடிகோளாக்கத்தையும் !!{derivation} முறைமையின் விவரிப்பையும்
கொடுத்தால், $!G_\Gamma$ ஐ உண்மையாகவே எழுதுவதற்கான ஒரு முறையை நிறுவல்
தருகிறது.

கோடலின் நிறுவலிலுள்ள கட்டமைப்புக்கு, $\Lang L_A$ இன் சொற்கள் மற்றும்
!!{formula}s மீதான பண்புகளையும் செயல்களையும் அதே~$\Lang L_A$ இல்
வெளிப்படுத்தும் வழி தேவை. ``$!A$ ஓர் !!a{sentence},'' ``$\delta$ என்பது
$!A$ இன் !!a{derivation},'' என்பன போன்ற பண்புகளும்,
$\Subst{!A}{t}{x}$ போன்ற செயல்களும் இவற்றில் அடங்கும். இந்த வழி
(a)~குறிகளையும் அவற்றின் தொடர்களையும் (சொற்கள், !!{formula}s,
!!{derivation}s போன்றவை இத்தகைய தொடர்களே) இயல் எண்களாக ($\Lang L_A$
பேசக்கூடியவை இவையே) ``குறியீடாக்குவதன்'' மூலம் இப்பண்புகளையும்
தொடர்புகளையும் வெளிப்படுத்த வேண்டும். மேலும், (b)~தொடர்புடைய உண்மைகளை
$\Gamma$ நிறுவும் வகையில் இதைச் செய்ய வேண்டும்; ஆகவே இப்பண்புகள் இயல்
எண்களின் !!{decidable} பண்புகளால் குறியீடாக்கப்படுகின்றன என்றும், செயல்கள்
இயல் எண்கள் மீதான கணிக்கத்தக்க சார்புகளுக்கு ஒத்துள்ளன என்றும் காட்ட
வேண்டும். இது ``தொடரமைப்பின் எண்கணிதமாக்கம்'' எனப்படுகிறது.

ஆனால் தொடரமைப்பை எவ்வாறு எண்கணிதமாக்கலாம் என்று ஆராய்வதற்கு முன்,
$\Gamma$ ``போதிய வலிமையுள்ளது,'' அதாவது எல்லாக் கணிக்கத்தக்க சார்புகளையும்
!!{decidable} தொடர்புகளையும் !!{represents} செய்கிறது என்ற நிபந்தனையை
ஆராய்வோம். இதற்கு ``கணிக்கத்தக்கது'' என்பதற்குத் துல்லியமான வரையறை தர
வேண்டும். டியூரிங் பொறிகளின் மாதிரியின் வழியாகவோ, ஏதாவது ஒரு பொதுப்பயன்
நிரலாக்க மொழியில் நிரல்களால் கணிக்கப்படும் சார்புகள் என்றோ இதைப் பல
வழிகளில் செய்யலாம். ஆனால் இச்சார்புகளையும் தொடர்புகளையும்~$\Lang L_A$
மொழியிலுள்ள ஒரு கோட்பாட்டில் !!{represents} செய்வதே நமது நோக்கம் என்பதால்,
எண்ணியல் சார்புகளின் கணிப்புத்தன்மைக்கான எளிய வரையறையைத் தேர்வது சிறந்தது.
இது \emph{மீள்வரையறைச் சார்பு} என்ற கருத்தாகும். ஆகவே முதலில்
மீள்வரையறைச் சார்புகளை விவாதிப்போம். பின்னர் எல்லா மீள்வரையறைச்
சார்புகளையும் தொடர்புகளையும் $\Th{Q}$ ஏற்கெனவே !!{represents} செய்கிறது
என்று காட்டுவோம். இவை ``போதிய வலிமையுள்ள'' கோட்பாடுகளுக்கான எடுத்துக்காட்டுகள்
என்பதை நிறுவியிருப்பதால், $\Th{Q}$, $\Th{PA}$ போன்ற குறிப்பிட்ட
கோட்பாடுகளுக்கு முழுமையின்மைத் தேற்றத்தைப் பயன்படுத்த இது உதவும்.

தொடரமைப்பின் எண்கணிதமாக்கத்தின் இறுதி விளைவு~$\OProv[\Gamma](x)$ என்ற
!!a{formula} ஆகும். !!{formula}s ஐ எண்களாகக் குறியீடாக்குவதன் வழியாக,
$\Gamma$ வின் அடிகோள்களிலிருந்து வருவிக்கத்தக்கதன்மையை அது வெளிப்படுத்துகிறது.
குறிப்பாக, $!A$ ஆனது~$n$ என்ற எண்ணால் குறியீடாக்கப்பட்டும் $\Gamma \Proves
!A$ என்றும் இருந்தால், $\Gamma \Proves \Prov[\Gamma](\num{n})$.
$\Gamma$ வுக்கான இந்த ``வருவிக்கத்தன்மைப் பயனிலை,'' ஒரு குறிப்பிட்ட
பொருளில் $\Gamma$ வின் முரணின்மையையும்~$\Lang L_A$ இன் !!a{sentence} ஆக
வெளிப்படுத்த உதவுகிறது: $\Gamma$ வுக்கான ``முரணின்மைக் கூற்று'' என்பது
$\lnot \Prov[\Gamma](\num{n})$ என்ற !!{sentence} என்க; இங்கு~$n$ என்பது
ஒரு முரண்பாட்டின், எடுத்துக்காட்டாக~$\lfalse$ இன், குறியீடாக எடுத்துக்கொள்ளப்படுகிறது.
முரணற்ற !!{axiomatizable} கோட்பாடுகள் தமது சொந்த முரணின்மைக் கூற்றுகளையும்
நிறுவுவதில்லை என்று இரண்டாம் முழுமையின்மைத் தேற்றம் கூறுகிறது. இத்தேற்றம்
பொருந்துவதற்குத் தேவையான நிபந்தனைகள், எல்லாக் கணிக்கத்தக்க சார்புகளையும்
!!{decidable} தொடர்புகளையும் கோட்பாடு பிரதிநிதித்துவப்படுத்துவதைவிடச் சற்று
வலுவானவை; ஆனால் $\Th{PA}$ அவற்றை நிறைவுசெய்கிறது என்று காட்டுவோம்.

\end{document}
